\documentclass{article}
\usepackage[T1]{fontenc}
\usepackage{lmodern}
\usepackage{graphicx}
\usepackage{amsmath,amssymb}
\usepackage[a4paper,margin=2cm]{geometry}
\usepackage[absolute,overlay]{textpos}
\setlength{\TPHorizModule}{1mm}
\setlength{\TPVertModule}{1mm}
\usepackage[indonesian]{babel}
\usepackage{hyperref}
\setlength{\emergencystretch}{2em}
% unit: Halaman 49

\begin{document}

% === Halaman 49 (llm) ===

\section*{Keamanan RC4}

\begin{itemize}
    \item RC4 memiliki kelemahan, yaitu \textit{byte-byte} pada \textit{keystream} awal pembangkitan memiliki korelasi yang tinggi dengan beberapa \textit{byte} awal kunci
    \item Oleh karena itu direkomendasikan membuang 256 sampai 512 \textit{byte keystream} awal
    \item RC4 adalah \textit{cipher} alir, maka ia tidak kuat terhadap serangan seperti \textit{flip-bit attack} maupun serangan-serangan \textit{stream attack} lainnya.
    \item Saat ini keamanan RC4 sudah berhasil dipecahkan dalam hitungan jam atau hari.
    \item Pada bulan Februari 2015, RC4 dilarang penggunaannya di dalam \textit{Transport Layer Security} (TLS) seperti disebutkan di dalam RFC 7465 (\url{https://tools.ietf.org/html/rfc7465}).
    \item Beberapa varian RC4 telah dibuat untuk mengatasi kelemahannya, yaitu Spritz, RC4A, VMPC, dan RC4+.
\end{itemize}

\end{document}
